% Generated draft: requires mathematical and source review.
\documentclass[11pt]{article}
\usepackage[margin=25mm]{geometry}
\usepackage{fontspec}
\setmainfont{Latin Modern Roman}
\usepackage{amsmath,amssymb,mathtools,unicode-math}
\setmathfont{Latin Modern Math}
\usepackage{enumitem,needspace,xurl,hyperref}
\setlength{\parindent}{0pt}
\setlength{\parskip}{5pt}
\setlength{\emergencystretch}{3em}
\newcommand{\N}{\mathbb{N}}
\newcommand{\Z}{\mathbb{Z}}
\newcommand{\Q}{\mathbb{Q}}
\newcommand{\R}{\mathbb{R}}
\newcommand{\C}{\mathbb{C}}
\newcommand{\F}{\mathbb{F}}
\newcommand{\A}{\mathbb{A}}
\newcommand{\PP}{\mathbb{P}}
\newcommand{\E}{\mathbb{E}}
\newcommand{\eps}{\varepsilon}
\newcommand{\dd}{\,\mathrm{d}}
\begin{document}
\section*{UnsolvedMath ID: 20003348}
\textbf{Source problem code:} catalogue entry
\subsection{Definitions and mathematical statement}
B. Weiss's Question 49 asks for cones $D\subset A$ of dimension at least three with no non-obvious divergent trajectories. The preceding definitions of $A$, ambient homogeneous space, cone, divergence and obviousness are not supplied locally. The words ``no non-obvious'' permit either only obvious divergent orbits or no divergent orbits, but these are different kinds of examples. The following sentence defining a radius-$T$ ball $A_T$ and Haar measure $\lambda$ belongs to the next question; the page header is layout.

The aid's proposed example uses $A\subset SL(4,\R)$, the positive determinant-one diagonal group, acting on unimodular lattices. Write
\[
 a(t)=\operatorname{diag}(e^{t_1},e^{t_2},e^{t_3},e^{t_4}),
 \qquad\sum_it_i=0.
\]
For the identity lattice it reports
\[
 \lambda_1(a(t)\Z^4)
 \leq\exp\!\left(-\frac{\|t\|_2}{\sqrt{12}}\right),
\]
where $\lambda_1$ is the Euclidean shortest-vector length. The four fixed rational vectors $e_1,e_2,e_3,e_4$ supply witnesses through their lengths $e^{t_i}$. This gives divergence as $\|t\|_2\to\infty$ along any proposed closed cone and realizes the aid's finite-witness obviousness at that lattice. It does not classify other lattice orbits.

The summary describes a specific closed three-dimensional weak fundamental-weight cone $D_4$, but does not give its inequalities. It says Datta--Tamam Theorem 1.16 would imply that every divergent $D_4$ orbit is obvious. The source/version comparison is unresolved: the exact v2 preprint is reported to print \emph{strict} inequalities in its introduction. No theorem statement, cone formula, divergence definition or preprint URL is supplied, so a weak-cone application is not asserted unconditionally and the missing inequalities are not guessed.

The aid's accessible-boundary lemma explains one definition issue. If a nonclosed strict cone $C$ has a finite boundary point $t_0\notin C$ approached by $t_j\in C$, the sequence escapes every compact subset of $C$ in its intrinsic topology, while $a(t_j)x\to a(t_0)x$. Therefore the orbit map from that nonclosed parameter space is not proper for any $x$. In this \emph{intrinsic properness} sense it has no divergent orbit. This is not the same as the asymptotic convention testing only $\|t\|\to\infty$ in the ambient Lie algebra. A sequence approaching a finite missing boundary point does not escape compact subsets of the ambient group.

Thus the closed cone has a reported explicit obvious divergent orbit, while the strict-cone/properness reading can be vacuous. Matching the cone closure and the theorem's actual divergence convention is essential before concluding that the candidate solves Question 49. The aid keeps the result partial and does not claim a complete universal no-non-obvious theorem for the closed candidate.
\subsection{Short English statement}
The aid gives an explicit shortest-vector bound for the identity lattice and a proposed closed three-dimensional cone. The all-orbits conclusion depends on an unresolved strict-versus-weak cone comparison. Properness for a nonclosed parameter space also differs from escape only at ambient infinity.
\subsection{Sources}
\begin{itemize}
\item Local UnsolvedMath registry, permanent ID 20003348, source code catalogue entry; \texttt{aim-other-notes.json}, notes[154], source tag question.
\item AIM, \emph{Emerging applications of measure rigidity}, Question 49, B. Weiss, printed page 20/PDF page 21 according to the supplied verification: \url{https://aimath.org/WWN/measrigid/measrigid.pdf}.
\item Datta--Tamam, Theorem 1.16 and the introduction of its v2 preprint, as reported by the aid; no full title, URL or theorem text is supplied.
\item Machine-generated research attempt 1: partial result, partially solved, rigor labelled ``proved,'' low novelty confidence. Not peer reviewed; independently verify before citation. Registry updated 14 August 2026.
\end{itemize}
\end{document}
